% Daftar Lampiran (dibuat manual dengan \pageref agar nomor halaman otomatis)
\chapter*{DAFTAR LAMPIRAN}
\addcontentsline{toc}{chapter}{DAFTAR LAMPIRAN}

\begingroup
\setlength{\parskip}{0.5em}
\noindent\textbf{Lampiran A.\hspace{0.4em}Kebutuhan Fungsional dan Non-Fungsional Sistem} \dotfill \pageref{lamp:fr}\par
\noindent\textbf{Lampiran B.\hspace{0.4em}Skema Basis Data} \dotfill \pageref{lamp:db}\par
\noindent\textbf{Lampiran C.\hspace{0.4em}Kamus Atribut Node dan Properti Temporal Knowledge Graph} \dotfill \pageref{lamp:kg}\par
\noindent\textbf{Lampiran D.\hspace{0.4em}Spesifikasi Antarmuka REST API dan Streaming} \dotfill \pageref{lamp:api}\par
\noindent\textbf{Lampiran E.\hspace{0.4em}Topologi Pipeline dan State Machine} \dotfill \pageref{lamp:pipeline}\par
\noindent\textbf{Lampiran F.\hspace{0.4em}Instrumen PHQ-9} \dotfill \pageref{lamp:phq9}\par
\noindent\textbf{Lampiran G.\hspace{0.4em}Instrumen Pengujian Pengguna} \dotfill \pageref{lamp:userstudy}\par
\noindent\textbf{Lampiran H.\hspace{0.4em}Skema dan Distribusi Korpus Linguistik} \dotfill \pageref{lamp:korpus}\par
\noindent\textbf{Lampiran I.\hspace{0.4em}Skema Audit Trail Guardrail} \dotfill \pageref{lamp:audit}\par
\noindent\textbf{Lampiran J.\hspace{0.4em}Visualisasi Knowledge Graph Pengguna} \dotfill \pageref{lamp:kg-visual}\par
\noindent\textbf{Lampiran K.\hspace{0.4em}Transkrip Lengkap Studi Kasus Evaluation Pipeline} \dotfill \pageref{lamp:eval-transcripts}\par
\noindent\textbf{Lampiran L.\hspace{0.4em}Contoh Prompt Sistem pada Pipeline Agentik} \dotfill \pageref{lamp:agentic-prompts}\par
\noindent\textbf{Lampiran M.\hspace{0.4em}Evaluasi Kritis Independen terhadap Studi Kasus Ilustratif} \dotfill \pageref{lamp:eval-critical}\par
\endgroup
